\documentclass[a4paper,11pt]{article}
\usepackage[utf8]{inputenc}
\usepackage[T1]{fontenc}
\usepackage[ngerman]{babel}
\usepackage{amsmath,amssymb}
\usepackage{xcolor}
\usepackage{tikz}
\usepackage{array}
\usepackage[a4paper,margin=2.2cm]{geometry}

\definecolor{bgdark}{HTML}{1E1E1E}
\definecolor{fgtext}{HTML}{E6E6E6}
\definecolor{gridcol}{HTML}{4A4A4A}
\definecolor{accent}{HTML}{6CB6FF}
\definecolor{leicht}{HTML}{7BD88F}
\definecolor{mittel}{HTML}{E6C07B}
\definecolor{schwer}{HTML}{E06C75}
\pagecolor{bgdark}
\color{fgtext}

\newcommand{\karo}[1]{\par\noindent
  \begin{tikzpicture}
    \draw[step=5mm, gridcol, very thin] (0,0) grid (\linewidth,#1);
  \end{tikzpicture}\par}
\newcommand{\aufgabe}[4]{\vspace{0.6em}\par\noindent
  \textbf{\large Aufgabe #1}\quad #2 \hfill {\color{#3}\textbf{[#4]}}\par\smallskip}
\setlength{\parindent}{0pt}

\begin{document}

\begin{center}
  {\color{accent}\Huge\bfseries Analysis II}\\[0.3em]
  {\Large Übungsaufgaben 03 — Integralrechnung}\\[0.8em]
\end{center}

\noindent\textbf{Name:} \rule{6cm}{0.4pt} \hfill \textbf{Datum:} \rule{3.5cm}{0.4pt}
\vspace{0.4em}
{\color{gridcol}\hrule height 1pt}
\vspace{0.3em}

\small\textit{Hilfsmittel: partielle Integration
$\int_a^b f'(x)g(x)\,dx = \big[f(x)g(x)\big]_a^b - \int_a^b f(x)g'(x)\,dx$;
Substitution $\int_a^b f(g(x))\,g'(x)\,dx = \int_{g(a)}^{g(b)} f(u)\,du$.}
\normalsize
\vspace{0.5em}

%==================================================================
\aufgabe{1}{Bestimmtes Integral (Potenzregel)}{leicht}{leicht}
Berechne
\[
  \int_{1}^{2} \left( 3x^2 - \frac{1}{x^2} \right) dx.
\]
\karo{5cm}

%==================================================================
\aufgabe{2}{Partielle Integration}{mittel}{mittel}
Berechne das unbestimmte Integral
\[
  \int x\,e^{2x}\,dx.
\]
\karo{5.5cm}

%==================================================================
\aufgabe{3}{Partielle Integration mit Logarithmus}{mittel}{mittel}
Berechne das unbestimmte Integral
\[
  \int x\,\ln(x)\,dx.
\]
\karo{5.5cm}

%==================================================================
\aufgabe{4}{Substitution}{mittel}{mittel}
Berechne das unbestimmte Integral
\[
  \int \frac{x^2}{x^3+1}\,dx.
\]
\karo{6cm}

%==================================================================
\aufgabe{5}{Trigonometrisches Integral mit Umformung}{schwer}{schwer}
Berechne das unbestimmte Integral
\[
  \int \left( \cos^3(x) + \frac{\sin^3(x)}{\tan(x)} \right) dx.
\]
\textit{Hinweis: Forme den Integranden zuerst geschickt um.}
\karo{6cm}

%==================================================================
\aufgabe{6}{Substitution mit Logarithmus}{mittel}{mittel}
Berechne das unbestimmte Integral
\[
  \int \frac{\ln(x)}{x}\,dx.
\]
\karo{5cm}

%==================================================================
\aufgabe{7}{Doppelintegral (separierbar)}{leicht}{leicht}
Berechne
\[
  \int_{0}^{2}\!\int_{0}^{1} x^2 y \;dx\,dy.
\]
\karo{5.5cm}

%==================================================================
\aufgabe{8}{Doppelintegral über eine Summe}{mittel}{mittel}
Berechne
\[
  \int_{0}^{1}\!\int_{0}^{2} (x + y) \;dx\,dy.
\]
\karo{5.5cm}

%==================================================================
\aufgabe{9}{Doppelintegral über ein Skalarfeld}{schwer}{schwer}
Berechne
\[
  \int_{0}^{1}\!\int_{0}^{1} (x^2 + y^2) \;dx\,dy.
\]
\karo{6cm}

%==================================================================
\aufgabe{10}{Zweifache partielle Integration}{schwer}{schwer}
Berechne das unbestimmte Integral
\[
  \int x^2\,e^{x}\,dx.
\]
\karo{6.5cm}

\end{document}
